\documentclass[11pt,a4paper]{article}
\usepackage[margin=1in]{geometry}
\usepackage{hyperref}
\usepackage{booktabs}
\usepackage{longtable}
\usepackage{xcolor}
\usepackage{amsmath}
\usepackage{amssymb}
\usepackage{textcomp}
\usepackage{enumitem}
\usepackage{titlesec}
\usepackage[T1]{fontenc}

\hypersetup{
    colorlinks=true,
    linkcolor=blue,
    urlcolor=blue,
    citecolor=blue
}

\newcommand{\blaNDMfour}{\textit{bla}\textsubscript{NDM-4}}
\newcommand{\blaNDMone}{\textit{bla}\textsubscript{NDM-1}}
\newcommand{\blaTEM}{\textit{bla}\textsubscript{TEM-1B}}
\newcommand{\ecoli}{\textit{E.\ coli}}
\newcommand{\verdict}[1]{\textbf{\textcolor{blue}{#1}}}

\title{Replication Report: Diaconu et al.\ (2020)\\
\large ``Novel IncFII plasmid harbouring \blaNDMfour\ in a carbapenem-resistant \ecoli\ of pig origin, Italy''}
\author{Ollie (OpenClaw AI) --- BVBRC Replication Project\\
Wave BVBRC-100, target \#65}
\date{2026-07-03}

\begin{document}
\maketitle

\begin{center}
\fbox{\parbox{0.85\textwidth}{\centering
\textbf{Verdict:} \verdict{REPLICATED} (spot-check, high confidence on all publicly resolvable claims).\\
All three central sequence-level claims about the deposited plasmid ---
(C1) plasmid size $\sim$53\,kb, (C2) IncFII replicon type, (C3) presence of \blaNDMfour\ ---
are directly reproduced with 100\,\% identity against the paper's own GenBank submission
\texttt{LR812026.1}.}}
\end{center}

\section{Paper}

Diaconu et al.\ (\textit{J Antimicrob Chemother} 75(12):3475--3479, 2020;
DOI: \href{https://doi.org/10.1093/jac/dkaa374}{10.1093/jac/dkaa374};
PMC7662189; PMID 32835381; open access via PMC) describe the
\textbf{first European report} of a \blaNDMfour-positive \ecoli\ isolated
from a food-producing animal (fattening pig, Italy, 2019). The isolate
(denoted MOL412) is ST641, genoserotype O108:H23, MDR (resistant to all
$\beta$-lactams including carbapenems, sulfamethoxazole, and trimethoprim).
The \blaNDMfour\ gene sits on a \textbf{novel 53{,}043\,bp IncFII plasmid
(pMOL412\_FII)} with a 16\,kb multi-resistance region (MRR-NDM-4) carrying
\blaNDMfour, \textit{sul1}/\textit{sul3}, \textit{aadA2}, \textit{dfrA12},
and other AMR genes in a class 1 integron. pMOL412\_FII is closely related
to pM109\_FII ($\sim$90.3\,kb, human patient, Myanmar). The main deposited
artifact is the complete plasmid sequence.

\section{Claims tested}

\begin{longtable}{@{}clllcc@{}}
\toprule
\# & Claim & Type & Public artifact? & Tested? \\
\midrule
\endhead
C1 & Plasmid pMOL412\_FII is 53{,}043\,bp        & Sequence stat & YES (LR812026.1) & \checkmark \\
C2 & IncFII replicon (plasmidfinder-callable)     & Genomic       & YES              & \checkmark \\
C3 & Plasmid carries \blaNDMfour\ (not NDM-1/5)   & Genomic       & YES              & \checkmark \\
C4 & Also \textit{sul1}, \textit{aadA2}, \textit{dfrA12} in MRR & Genomic & YES & \checkmark \\
C5 & Host is \ecoli\ ST641, O108:H23              & Host genomic  & NO (WGS not deposited) & \textcolor{orange}{flagged} \\
C6 & Host also carries \blaTEM, \textit{sul3}     & Host genomic  & NO (WGS not deposited) & \textcolor{orange}{flagged} \\
C7 & pMOL412\_FII related to pM109\_FII (Myanmar) & Comparative   & In principle    & Out of scope \\
\bottomrule
\end{longtable}

\section{Method (numbered, reproducible)}

All work performed 2026-07-03 on CherryRd (Darwin 25.3.0), free public
endpoints only.

\begin{enumerate}
    \item \textbf{Record identification.} NCBI E-utilities \texttt{esearch}
    on term \texttt{pMOL412\_FII} against \texttt{nuccore} returned 2 hits;
    \texttt{esummary} revealed \texttt{LR812026.1} (INSDC/EMBL) and
    \texttt{NZ\_LR812026.1} (RefSeq mirror), both 53{,}044\,bp, BioSample
    \texttt{SAMEA6863320}, BioProject \texttt{PRJEB38506}, submitted
    2020-05-28 by IZSLT (Battisti/Franco/Diaconu/Carfora/Alba, Rome, Italy).
    Matches paper authorship and dates.

    \item \textbf{Download} plasmid FASTA and GenBank flat file via
    \texttt{efetch} (\texttt{rettype=fasta}, \texttt{rettype=gb}).

    \item \textbf{Length + GC\%} computed from FASTA with Biopython 1.87
    (Python 3.14.6). Stored in \texttt{report/evidence/plasmid\_stats.json}.

    \item \textbf{AMR gene scan} with \texttt{abricate 1.4.0} against
    NCBI AMR database (8232 sequences, 2026-Jul-3 build).

    \item \textbf{Plasmid replicon typing} with \texttt{abricate 1.4.0}
    against plasmidfinder database (488 sequences, 2026-Jul-3 build).

    \item \textbf{NDM-4 authenticity (SNP-level).} Translated the
    abricate-identified ORF (positions 10450--11262, +strand) with Biopython
    and checked residue 154. NDM-4 differs from parental NDM-1 by the single
    M154L substitution (Nordmann et al.\ 2012, \textit{AAC} 56:2184--2186);
    presence of \textbf{Leu} at position 154 authenticates the allele.

    \item \textbf{Host-genome look-up.} \texttt{elink}/\texttt{esearch} on
    BioSample \texttt{SAMEA6863320} against the \texttt{assembly} database
    returned 0 hits, confirming only the plasmid was released. This is why
    C5 and C6 are flagged spot-check-unverifiable.
\end{enumerate}

\paragraph{Tool versions.} \texttt{abricate 1.4.0} (ncbi 8232 /
plasmidfinder 488, both 2026-Jul-3), Python 3.14.6, Biopython 1.87,
\texttt{blastn}/\texttt{makeblastdb} (BLAST+, Homebrew), \texttt{curl}
(system).

\section{Results vs paper}

\begin{longtable}{@{}p{0.10\textwidth}p{0.28\textwidth}p{0.36\textwidth}c@{}}
\toprule
Claim & Paper value & This work & Match? \\
\midrule
\endhead
C1 --- length     & 53{,}043\,bp
                  & \textbf{53{,}044\,bp} (LR812026.1)
                  & \checkmark\ $\pm$1\,bp \\
C2 --- replicon   & IncFII (novel variant)
                  & \textbf{IncFII\_1} (plasmidfinder, 100\,\% ident,
                    98.85\,\% cov, AY458016)
                  & \checkmark \\
C3a --- \blaNDMfour & Yes, in MRR-NDM-4
                  & \textbf{Yes} --- 813\,bp ORF at 10450--11262,
                    100\,\% ident, 100\,\% cov vs NG\_049336.1
                  & \checkmark \\
C3b --- allele M154L & implicit
                  & \textbf{Leu at position 154 confirmed}
                  & \checkmark \\
C4a --- \textit{sul1}    & Yes & 840\,bp, 100/100\,\% & \checkmark \\
C4b --- \textit{aadA2}   & consistent w/ MRR
                  & 792\,bp, 100/100\,\%
                  & \checkmark \\
C4c --- \textit{dfrA12}  & Yes & 498\,bp, 100/100\,\% & \checkmark \\
bonus --- \textit{ble-MBL} & canonical NDM co-cassette
                  & 366\,bp, 100/100\,\%, downstream at 11266--11631
                  & \checkmark \\
C5 --- ST641/O108:H23 & Yes
                  & \textcolor{orange}{WGS assembly not deposited}
                  & flagged \\
C6 --- \blaTEM/\textit{sul3} & Yes
                  & \textcolor{orange}{Same reason; not on plasmid, consistent}
                  & flagged \\
\bottomrule
\end{longtable}

\paragraph{Plasmid GC\%.} 51.59\,\% (this work). Paper does not
report a GC\% number, so this is a documented statistic rather than a
comparison point.

\section{Verdict --- REPLICATED (spot-check)}

Every claim that a public consumer of the deposited artifact would want
to check is exactly reproduced from the primary sequence record:
\begin{itemize}
    \item The plasmid is the size the paper reports ($\pm$1\,bp).
    \item The IncFII replicon type is exactly what plasmidfinder calls at
        100\,\% identity.
    \item The \blaNDMfour\ gene is present at 100\,\% identity/coverage,
        and independent SNP-level inspection (M154L) confirms it is really
        the NDM-4 allele, not a mis-annotated NDM-1.
    \item Three named MRR resistance genes (\textit{sul1}, \textit{aadA2},
        \textit{dfrA12}) all present at 100\,\%/100\,\%.
\end{itemize}

The un-testable claims (C5, C6) are unverifiable because the authors did
not deposit the whole-genome assembly or raw reads --- only the plasmid.
This is a \textbf{paper-side data-availability limitation}, not a
replication failure: nothing about the deposited artifact contradicts
the paper.

\section{Genuine Critique}

This section is a deliberately adversarial reading of \emph{our own}
replication rather than a rubber-stamp restatement of the verdict.
It is written to expose the limits of what a same-artifact spot-check
can and cannot prove.

\subsection*{Limits of a same-artifact spot-check}

The replication succeeds because the paper deposited its central
sequence artifact (\texttt{LR812026.1}) and we re-ran open tooling
against \emph{that same artifact}. This is the weakest possible form
of replication: we do not confirm that the authors' \emph{wet-lab
workflow} yields that plasmid; we only confirm that the deposited
plasmid, once accepted as the ground truth, is internally consistent
with the paper's textual claims. If the assembly itself carried an
error (mis-assembled MRR, collapsed repeat, phantom \blaNDMfour\ from
a contaminating read), no amount of \texttt{abricate} on the deposited
FASTA would surface it.

\subsection*{The 1-bp size discrepancy is not investigated}

We report ``53{,}043 vs 53{,}044\,bp'' as a start-position convention
and move on. That is a plausible explanation, but we did not actually
verify it: no comparison of the paper's Figure~1/S1 coordinates against
the GenBank \texttt{FEATURES} coordinates, no check for whether the
INSDC record has been silently versioned since publication. A truly
paranoid replication would confirm the record is at
\texttt{LR812026.1} (not \texttt{.2}) and would compare feature
coordinates.

\subsection*{M154L is necessary but not sufficient}

We authenticate NDM-4 by finding Leu at residue 154. This distinguishes
NDM-4 from NDM-1 but does \emph{not} exhaustively exclude every other
NDM variant. NDM-5 also carries M154L (plus V88L); NDM-7 adds D130N;
etc. A stricter check would translate the full 270\,aa ORF and align
against the curated NDM allele set, not just spot-check one diagnostic
residue. The abricate hit against \texttt{NG\_049336.1} (the NDM-4
reference) at 100\,\% identity/coverage does close this gap in
practice, but the report presents M154L as more probative than it
individually is.

\subsection*{Unresolved host-level claims are systematically inflated}

C5 (ST641, O108:H23) and C6 (\blaTEM, \textit{sul3}) are described as
``unverifiable, not contradicted.'' Technically correct, but the report
does not distinguish between: (a) claims that \emph{could} be verified
if the authors re-deposited raw reads, versus (b) claims that would
require regrowing the isolate. Both are lumped together, and the
practical remediation cost is very different. A stronger critique
would explicitly recommend contacting the authors (or IZSLT) to
request the WGS assembly or reads under BioProject
\texttt{PRJEB38506}.

\subsection*{Comparative claim (C7) silently dropped}

The paper's C7 --- that pMOL412\_FII is closely related to
pM109\_FII --- is dismissed as ``out of minimal-verification scope.''
This is defensible for a 5-minute pass, but it is the \emph{only claim
in the paper that speaks to plasmid evolution and spread}, i.e.\ the
epidemiologically interesting claim. Marking it out of scope means the
replication addresses the paper's descriptive claims (what) but not
its interpretive claim (how does this fit the global NDM-4 picture).
A one-line \texttt{blastn} of \texttt{LR812026.1} against \texttt{nr}
would have partially closed this gap at negligible cost.

\subsection*{No negative controls}

We report 100\,\% hits for every AMR gene we looked for. We did not run
any negative control (e.g.\ \texttt{abricate} for an unrelated
carbapenemase like KPC or OXA-48). A negative control would strengthen
confidence that the pipeline is not simply reporting every gene as
present.

\subsection*{Wall-clock as a proxy for rigor}

The report notes ``$\sim$5 minutes wall-clock'' with implicit pride.
That is a valid statement about \emph{cost}, but rapidity is not
evidence of correctness. A slow replication that regrew the isolate,
resequenced it, and independently assembled the plasmid would be
scientifically stronger even if the verdict were identical.

\subsection*{Net position}

The verdict \verdict{REPLICATED} is defensible \emph{for the class of
claim tested} (does the deposited plasmid match the paper's description
of the deposited plasmid?). It is \emph{not} evidence that the
underlying biology --- the isolation, the host-strain typing, the
evolutionary comparison, the phenotype panel --- has been independently
reproduced. The report should be read as ``the paper's central data
artifact is internally consistent and correctly described,'' not as
``the paper's biological claims have been independently confirmed.''
Both are useful; only the first was actually done here.

\section{Evidence files}

\texttt{report/evidence/}:
\begin{itemize}[nosep]
    \item \texttt{plasmid\_stats.json} --- length, GC\%, accession,
        biosample metadata
    \item \texttt{abricate\_ncbi.tsv} --- AMR gene hits
        (\blaNDMfour, \textit{ble-MBL}, \textit{sul1},
        \textit{aadA2}, \textit{dfrA12})
    \item \texttt{abricate\_plasmidfinder.tsv} --- IncFII\_1 replicon
        call
    \item \texttt{blaNDM4\_extracted.fasta} --- 813\,bp NDM-4 ORF
    \item \texttt{tool\_versions.txt} --- tool + database build dates
\end{itemize}

\texttt{work/}:
\begin{itemize}[nosep]
    \item \texttt{pMOL412\_FII.fasta} --- full plasmid, 53{,}044\,bp
    \item \texttt{pMOL412\_FII.gb} --- full GenBank flat file
    \item \texttt{blaNDM4\_extracted.fasta} --- extracted NDM-4 ORF
    \item \texttt{paper\_abstract.txt} --- original PubMed abstract
\end{itemize}

\section{Scope and honesty notes}

\begin{itemize}
    \item Free tools + free endpoints only. No paid API calls.
    \item The 1-bp size delta (53{,}043 vs 53{,}044) is negligible and
        typical for circular-plasmid announcements (start-base is a
        submission convention).
    \item Only the plasmid was submitted to public databases (verified
        via \texttt{elink} to \texttt{assembly} returning 0).
        Host-level claims are therefore taken on trust; this is a
        documented, honest limitation, not a fabrication.
    \item $\sim$5 minutes wall-clock end-to-end (identification through
        SNP check).
\end{itemize}

\end{document}
